\documentclass{article}
\usepackage{geometry}
\geometry{margin=1.5cm, vmargin={0pt,1cm}}
\setlength{\topmargin}{-1cm}
\setlength{\headheight}{12.5pt}
\setlength{\paperheight}{29.7cm}
\setlength{\textheight}{25.3cm}

% useful packages.
\usepackage{fancyhdr}
\usepackage{layout}
\usepackage{tikz}
\usetikzlibrary{arrows.meta}

% import common macros
% % % % % % % % % % % % % % % % % % % % % % % % % % % % % % % %
% Common Macros for LaTeX
% % % % % % % % % % % % % % % % % % % % % % % % % % % % % % % %

% Useful packages
\usepackage{amsfonts}
\usepackage{amsmath}
\usepackage{amssymb}
\usepackage{amsthm}
\usepackage{enumerate}
\usepackage{graphicx}
\usepackage{multicol}
\usepackage[ruled,linesnumbered]{algorithm2e}

% Math operators and common symbols
\newcommand{\dif}{\mathrm{d}}
\newcommand{\innerProd}[1]{\left\langle #1 \right\rangle}
\newcommand{\difFrac}[2]{\frac{\dif #1}{\dif #2}}
\newcommand{\pdfFrac}[2]{\frac{\partial #1}{\partial #2}}
\newcommand{\T}{\mathrm{T}}
\newcommand{\norm}[1]{\left\| #1 \right\|}
\newcommand{\abs}[1]{\left| #1 \right|}

% Vector and Matrix shortcuts
\newcommand{\vecTwo}[2]{
  \begin{bmatrix}
    #1 \\
    #2
\end{bmatrix}}
\newcommand{\vecThree}[3]{
  \begin{bmatrix}
    #1 \\
    #2 \\
    #3
\end{bmatrix}}
\newcommand{\vecTwoH}[2]{
  \begin{bmatrix}
    #1 & #2
  \end{bmatrix}
}
\newcommand{\vecThreeH}[3]{
  \begin{bmatrix}
    #1 & #2 & #3
  \end{bmatrix}
}

% Common fields
\newcommand{\R}{\mathbb{R}}
\newcommand{\C}{\mathbb{C}}
\newcommand{\Z}{\mathbb{Z}}
\newcommand{\N}{\mathbb{N}}

% Solutions and proofs
\newenvironment{solution}{
\begin{proof}[解]}{
\end{proof}}

% Utility to get environment variables (requires lualatex)
\newcommand{\getEnv}[2]{\directlua{tex.print(os.getenv("#1") or "#2")}}


\newcommand{\Res}{\operatorname{Res}}
\newcommand{\PV}{\operatorname{P.V.}}
\newcommand{\Chat}{\widehat{\C}}

% Name and uID
\newcommand{\name}{\getEnv{NAME}{NAME}}
\newcommand{\uid}{\getEnv{UID}{UID}}
\newcommand{\course}{Complex Analysis}
\newcommand{\hwNum}{12}

\begin{document}

\pagestyle{fancy}
\fancyhead{}
\lhead{\zhstr{\name} (\uid)}
\chead{\course\ \#\hwNum}
\rhead{\today}

\section{Corollary 4.134}

The residue of a simple pole $c$ of $f$ is
\[
  \Res_{z=c} f(z) = \lim_{z\to c}(z-c)f(z).
\]

\begin{solution}
  Since $c$ is a simple pole, by the definition of pole order there exists
  $r>0$ such that
  \[
    f(z)=\sum_{n=-1}^{\infty} a_n(z-c)^n,\qquad 0<|z-c|<r.
  \]
  The residue is the Laurent coefficient $a_{-1}$ by Definition~4.131.
  Multiplying by $z-c$ gives
  \[
    (z-c)f(z)=a_{-1}+\sum_{n=0}^{\infty} a_n(z-c)^{n+1}.
  \]
  Letting $z\to c$ yields
  \[
    \lim_{z\to c}(z-c)f(z)=a_{-1}=\Res_{z=c}f(z).
  \]
  Equivalently, this is Lemma~4.133 with pole order $m=1$.
\end{solution}

\section{Exercise 4.153}

Prove that
\[
  \int_{-\infty}^{+\infty}\frac{\dd x}{(1+x^2)^{n+1}}
  =
  \frac{1\cdot 3\cdot 5\cdots(2n-1)}
       {2\cdot 4\cdot 6\cdots(2n)}\pi .
\]

\begin{solution}
  Let
  \[
    f(z)=\frac{1}{(1+z^2)^{n+1}}
    =\frac{1}{(z-i)^{n+1}(z+i)^{n+1}}.
  \]
  The hypotheses of Theorem~4.150 apply: the denominator has no real zeros
  and its degree is $2n+2$, at least two larger than the numerator degree.
  The only pole in the upper half plane is $z=i$, of order $n+1$. Hence
  \[
    \int_{-\infty}^{+\infty}\frac{\dd x}{(1+x^2)^{n+1}}
    =2\pi i\,\Res_{z=i} f(z).
  \]

  By Lemma~4.133, with
  \[
    \phi(z)=(z-i)^{n+1}f(z)=(z+i)^{-n-1},
  \]
  we have
  \[
    \Res_{z=i} f(z)
    =\frac{\phi^{(n)}(i)}{n!}.
  \]
  Since
  \[
    \frac{\dd^n}{\dd z^n}(z+i)^{-n-1}
    =(-1)^n\frac{(2n)!}{n!}(z+i)^{-2n-1},
  \]
  it follows that
  \[
    \Res_{z=i} f(z)
    =
    \frac{(-1)^n(2n)!}{(n!)^2}(2i)^{-2n-1}
    =
    \frac{(2n)!}{2^{2n+1}(n!)^2\,i}.
  \]
  Therefore
  \[
    \int_{-\infty}^{+\infty}\frac{\dd x}{(1+x^2)^{n+1}}
    =
    2\pi i\cdot
    \frac{(2n)!}{2^{2n+1}(n!)^2\,i}
    =
    \pi\frac{(2n)!}{2^{2n}(n!)^2}.
  \]
  Finally,
  \[
    \frac{(2n)!}{2^{2n}(n!)^2}
    =
    \frac{1\cdot 3\cdot 5\cdots(2n-1)}
         {2\cdot 4\cdot 6\cdots(2n)},
  \]
  which proves the desired formula.
\end{solution}

\section{Exercise 4.155}

Prove Lemma~4.154, i.e.,
\[
  \forall x\in\left[0,\frac{\pi}{2}\right],
  \qquad
  \frac{2}{\pi}x\le \sin x\le x.
\]

\begin{solution}
  Define
  \[
    h(x)=x-\sin x.
  \]
  Then $h(0)=0$ and
  \[
    h'(x)=1-\cos x\ge 0
    \qquad
    \left(0\le x\le \frac{\pi}{2}\right).
  \]
  Thus $h(x)\ge 0$, proving $\sin x\le x$.

  For the lower bound, define
  \[
    g(x)=\sin x-\frac{2}{\pi}x.
  \]
  Then $g(0)=g(\pi/2)=0$ and
  \[
    g''(x)=-\sin x\le 0
    \qquad
    \left(0\le x\le \frac{\pi}{2}\right).
  \]
  Hence $g$ is concave on $[0,\pi/2]$. A concave function lies above the
  chord joining two points on its graph, so $g(x)\ge 0$ on the interval.
  Therefore
  \[
    \sin x\ge \frac{2}{\pi}x.
  \]
  Combining the two inequalities proves Lemma~4.154.
\end{solution}

\section{Corollary 4.158}

For all $m>0$, prove that
\[
  \int_{-\infty}^{+\infty}\frac{\cos(mx)}{1+x^2}\,\dd x
  =\pi e^{-m}.
\]

\begin{solution}
  Consider
  \[
    f(z)=\frac{e^{imz}}{1+z^2}.
  \]
  The denominator has no real zeros, and $f$ has only one pole in the upper
  half plane, namely the simple pole $z=i$. By Theorem~4.157,
  \[
    \int_{-\infty}^{+\infty}\frac{e^{imx}}{1+x^2}\,\dd x
    =
    2\pi i\,\Res_{z=i}f(z).
  \]
  Using Corollary~4.136 for the simple pole,
  \[
    \Res_{z=i}f(z)
    =
    \frac{e^{imz}}{2z}\bigg|_{z=i}
    =
    \frac{e^{-m}}{2i}.
  \]
  Hence
  \[
    \int_{-\infty}^{+\infty}\frac{e^{imx}}{1+x^2}\,\dd x
    =
    2\pi i\cdot \frac{e^{-m}}{2i}
    =
    \pi e^{-m}.
  \]
  Taking real parts gives
  \[
    \int_{-\infty}^{+\infty}\frac{\cos(mx)}{1+x^2}\,\dd x
    =\pi e^{-m}.
  \]
  The imaginary part is zero, consistently with the integrand
  $\sin(mx)/(1+x^2)$ being odd.
\end{solution}

\section{Exercise 4.159}

Show that
\[
  \int_{-\infty}^{+\infty}
  \frac{x\cos x}{x^2-2x+10}\,\dd x
  =
  \frac{\pi}{3}e^{-3}(\cos 1-3\sin 1).
\]

\begin{solution}
  Let
  \[
    F(z)=\frac{z e^{iz}}{z^2-2z+10}.
  \]
  Since
  \[
    z^2-2z+10=(z-(1+3i))(z-(1-3i)),
  \]
  the only pole of $F$ in the upper half plane is the simple pole
  \[
    a=1+3i.
  \]
  The hypotheses of Theorem~4.157 are satisfied with $k=1$, so
  \[
    \int_{-\infty}^{+\infty}
    \frac{x e^{ix}}{x^2-2x+10}\,\dd x
    =
    2\pi i\,\Res_{z=a}F(z).
  \]
  By Corollary~4.136,
  \[
    \Res_{z=a}F(z)
    =
    \frac{a e^{ia}}{2a-2}
    =
    \frac{(1+3i)e^{i(1+3i)}}{6i}.
  \]
  Therefore
  \[
    \int_{-\infty}^{+\infty}
    \frac{x e^{ix}}{x^2-2x+10}\,\dd x
    =
    2\pi i\cdot
    \frac{(1+3i)e^{i(1+3i)}}{6i}
    =
    \frac{\pi}{3}(1+3i)e^{-3}e^{i}.
  \]
  Taking real parts yields
  \[
    \int_{-\infty}^{+\infty}
    \frac{x\cos x}{x^2-2x+10}\,\dd x
    =
    \frac{\pi}{3}e^{-3}
    \Re\bigl((1+3i)(\cos 1+i\sin 1)\bigr).
  \]
  Since
  \[
    \Re\bigl((1+3i)(\cos 1+i\sin 1)\bigr)
    =
    \cos 1-3\sin 1,
  \]
  the desired identity follows.
\end{solution}

\section{Definition 4.165}

In the 2026-06-21 complete handout, item~4.165 is the definition of open
sets in the extended complex plane $\Chat$:
\[
  D\subseteq \Chat \text{ is open in } \Chat
  \iff
  D\cap \C \text{ is open in } \C
  \text{ and }
  \infty\in D \Rightarrow
  \exists R>0,\ \{z\in\C:|z|>R\}\subseteq D.
\]

\begin{solution}
  This item is a definition rather than a proposition. Its content is that
  ordinary finite points keep their usual Euclidean neighborhoods, while a
  neighborhood of $\infty$ must contain all sufficiently large finite points.

  Equivalently, the basic neighborhoods of $\infty$ are
  \[
    U_R(\infty):=\{\infty\}\cup\{z\in\C:|z|>R\},
    \qquad R>0.
  \]
  Thus the condition in Definition~4.165 says exactly that if an open set
  contains $\infty$, then it contains one of these neighborhoods of $\infty$.
  This is the topology used in Definition~4.180 to classify singularities at
  infinity by the substitution $z\mapsto 1/z$.
\end{solution}

\section{Example 4.173}

Show that
\[
  f(z)=\cot(\pi z)=\frac{\cos(\pi z)}{\sin(\pi z)}
\]
is not meromorphic in $\Chat$.

\begin{solution}
  By Definition~4.170, to decide whether $f$ is meromorphic in $\Chat$ we must
  inspect
  \[
    \widehat f(z)=f(1/z)=\cot\frac{\pi}{z}
  \]
  near $z=0$, because $z=0$ corresponds to $\infty$.

  The finite poles of $\widehat f$ occur when
  \[
    \sin\frac{\pi}{z}=0
    \iff
    \frac{\pi}{z}=n\pi
    \iff
    z=\frac1n,
    \qquad n\in\Z\setminus\{0\}.
  \]
  Hence
  \[
    S(\widehat f)\supseteq
    \left\{\frac1n:n\in\Z\setminus\{0\}\right\}.
  \]
  This set accumulates at $0$. Since $0$ belongs to the domain used to test
  meromorphicity at $\infty$, the pole set is not discrete. This violates
  condition~(a) in Definition~4.169.

  Therefore $\widehat f$ is not meromorphic near $0$, and by
  Definition~4.170 the original function $\cot(\pi z)$ is not meromorphic in
  the extended complex plane $\Chat$.
\end{solution}

\section{Theorem 4.176}

Let $\Omega\subseteq\C$ be a nonempty open set. Then $H(\Omega)$ is an
integral domain if and only if $\Omega$ is a domain. In that case
$H(\Omega)=O(\Omega)$.

\begin{solution}
  By Lemma~4.174, $H(\Omega)$ is a commutative ring with unity. Hence the
  only additional property needed for $H(\Omega)$ to be an integral domain is
  the absence of divisors of zero.

  \medskip
  \noindent\textbf{Sufficiency.}
  Assume $\Omega$ is a domain, i.e. $\Omega$ is connected and open. Let
  $f,g\in H(\Omega)$ and suppose $fg=0$ on $\Omega$.
  If $f\equiv 0$, then there is nothing to prove. Otherwise, choose
  $a\in\Omega$ with $f(a)\ne 0$. By continuity of $f$, there is a small disk
  $U_r(a)\subseteq \Omega$ on which $f$ has no zeros. Since $fg=0$, we get
  $g=0$ on $U_r(a)$.

  The identity theorem (Theorem~4.35) now gives $g\equiv 0$ on the connected
  domain $\Omega$. Thus $fg=0$ implies $f=0$ or $g=0$, so there are no
  divisors of zero. Therefore $H(\Omega)$ is an integral domain.

  \medskip
  \noindent\textbf{Necessity.}
  Assume $\Omega$ is not connected. Then there exist nonempty disjoint open
  subsets $A,B\subseteq\Omega$ such that
  \[
    \Omega=A\cup B.
  \]
  Since $\Omega$ is open, $A$ and $B$ are open in $\C$ as well. Define
  \[
    f(z)=
    \begin{cases}
      1,& z\in A,\\
      0,& z\in B,
    \end{cases}
    \qquad
    g(z)=
    \begin{cases}
      0,& z\in A,\\
      1,& z\in B.
    \end{cases}
  \]
  Both $f$ and $g$ are locally constant, hence analytic on $\Omega$. They are
  both nonzero elements of $H(\Omega)$, but
  \[
    fg\equiv 0.
  \]
  Thus $H(\Omega)$ has divisors of zero, contradicting the definition of an
  integral domain. Hence $\Omega$ must be connected, i.e. a domain.

  Finally, when $\Omega$ is a domain, Notation~6 gives
  $H(\Omega)=O(\Omega)$, since both denote the analytic functions on the same
  domain.
\end{solution}

\section{Example 4.182}

Show that the entire function $f(z)=e^z$ has an essential singularity at
$\infty$.

\begin{solution}
  By Definition~4.180, the singularity of $f$ at $\infty$ is classified by
  the singularity of
  \[
    \widehat f(z)=f(1/z)=e^{1/z}
  \]
  at $z=0$.

  For $z\ne 0$,
  \[
    e^{1/z}
    =
    \sum_{n=0}^{\infty}\frac{1}{n!z^n}
    =
    1+\frac1z+\frac{1}{2!z^2}+\frac{1}{3!z^3}+\cdots .
  \]
  This Laurent expansion has infinitely many negative-power terms. By the
  Laurent classification theorem (Theorem~4.126), the singularity of
  $\widehat f$ at $0$ is essential, not removable and not a pole.
  Therefore $e^z$ has an essential singularity at $\infty$.
\end{solution}

\section{Corollary 4.187}

Any analytic function $f:\Chat\to\C$ is constant.

\begin{solution}
  Since $f$ is analytic on the extended complex plane, its restriction
  $f|_{\C}$ is entire. Also, $f$ is analytic at $\infty$, so the singularity
  of $f|_{\C}$ at $\infty$ is removable in the sense of Definition~4.180.

  By Theorem~4.184, an entire function has a removable singularity at
  $\infty$ if and only if it is constant. Applying this to $f|_{\C}$ gives
  that $f|_{\C}$ is constant. Since $f$ is continuous at $\infty$, the value
  at $\infty$ must equal the same constant. Hence $f$ is constant on $\Chat$.

  Equivalently, by Theorem~4.185, an analytic function on $\Chat$ is
  meromorphic and therefore rational; having no poles anywhere on $\Chat$
  forces it to be a constant rational function.
\end{solution}

\section{Theorem 4.189}

An entire function $f:\C\to\C$ is globally conformal if and only if
\[
  \exists a\in\C^\ast,\ \exists b\in\C
  \quad\text{such that}\quad
  f(z)=az+b.
\]

\begin{solution}
  \textbf{Sufficiency.}
  If $f(z)=az+b$ with $a\ne 0$, then $f$ is entire,
  \[
    f'(z)=a\ne 0,
  \]
  and it has the entire inverse
  \[
    f^{-1}(w)=\frac{w-b}{a}.
  \]
  Thus $f$ is globally conformal.

  \medskip
  \noindent\textbf{Necessity.}
  Assume $f$ is globally conformal. Then $f$ is an entire bijection and its
  inverse is continuous. We first show that
  \[
    \lim_{z\to\infty} f(z)=\infty.
  \]
  If not, there would exist a sequence $z_n$ with $|z_n|\to\infty$ but
  $\{f(z_n)\}$ bounded. Passing to a subsequence, assume
  $f(z_n)\to w\in\C$. Since $f^{-1}$ is continuous,
  \[
    z_n=f^{-1}(f(z_n))\to f^{-1}(w),
  \]
  contradicting $|z_n|\to\infty$.

  Therefore the singularity of $f$ at $\infty$ is a pole rather than a
  removable singularity or an essential singularity. By Theorem~4.184, $f$ is
  a nonconstant polynomial.

  It remains to show that this polynomial has degree $1$. If
  $\deg f=d\ge 2$, then $f'$ is a nonconstant polynomial of degree $d-1$.
  By the fundamental theorem of algebra, $f'$ has a zero. This contradicts
  conformality, since a conformal map must have nonzero derivative at every
  finite point. Hence $d=1$, and so
  \[
    f(z)=az+b
  \]
  for some $a\ne 0$ and $b\in\C$.
\end{solution}

\end{document}
